\answersfor{ch:submission}

\begin{solution}{ex:submission:1}
\begin{align*}
\text{span} &= 429{,}985-400{,}015 = 29{,}970\text{ m},\\
\Delta x &= 29{,}970/(1000-1) = 29{,}970/999 = 30.0\text{ m},\\
c &= 400{,}015-15 = 400{,}000 .
\end{align*}
Dividing by $W$: $29{,}970/1000=29.97$\,m. If the UL value is also taken as the image's left edge ($c=400{,}015$), the
last column's centre lands at
\begin{align*}
400{,}015+999.5\times29.97 = 400{,}015+29{,}955.015 = 429{,}970.015,
\end{align*}
against the attribute's 429,985: $14.985$\,m, half a pixel, off. The first column's centre,
$400{,}015+0.5\times29.97=400{,}029.985$, is off by the same amount: the whole grid slides half a pixel.
\end{solution}

\begin{solution}{ex:submission:2}
The centre is at $j=20.5$, $i=10.5$:
\begin{align*}
X &= 30\times20.5+545{,}007.75 = 615+545{,}007.75 = 545{,}622.75,\\
Y &= -30\times10.5+2{,}501{,}437.5 = -315+2{,}501{,}437.5 = 2{,}501{,}122.5,
\end{align*}
in EPSG:32642.
\end{solution}

\begin{solution}{ex:submission:3}
The omitted field takes the pydantic default \code{True}, so the vendable is on the 165-band grid. v0.2.0 could run on
it, but the worker checks the models in list order, so V2 comes first and \code{_check_band_grid} refuses it
(``reads native bands: run band\_filter\_apply with use\_common\_wavelength\_grid=false''): the action fails. Fix: a band filter with the field
\code{false}, which is what the UI's defaults give.
\end{solution}

\begin{solution}{ex:submission:4}
8-connected: the centre pixel touches all four corners, so all five pixels form \textbf{1} object. 4-connected: no two
pixels share an edge, so there are \textbf{5} objects.
\end{solution}

\begin{solution}{ex:submission:5}
$0.03\times10000=300\to300$. $-1.2\times10000=-12{,}000\to-9998$ (the lower clip, so it cannot be read as nodata).
$3.5\times10000=35{,}000\to32{,}767$ (int16's largest value).
\end{solution}

\begin{solution}{ex:submission:6}
A commit overwrites the prep's \file{anomaly_mask.tif}, so result~1 would be lost, and the export names the zip by the
action's short id, so two results from one action would not be told apart. Three results need three prep actions.
\end{solution}
